\section{一个可以手算的严格改进}
取 $2\times2$ 网格，四条边均为 $h=1/5$。这时每个单元的两个外边界面各贡献 $2/21$，算子为方形四点图的 Laplacian 加上 $(4/21)I$。考虑从左到右、从下到上编号的三种残差：
\[
 r_{\rm same}=\begin{pmatrix}1&1\\1&1\end{pmatrix},\qquad
 r_{\rm alt}=\begin{pmatrix}1&-1\\-1&1\end{pmatrix},\qquad
 r_{\rm one}=\begin{pmatrix}1&0\\0&0\end{pmatrix}.
\]
它们的最大绝对残差都是 $1$。前两者分别是图 Laplacian 对应于特征值 $0$ 和 $4$ 的向量，因此
\[
 \|G_{\rm w}r_{\rm same}\|_\infty=\frac{21}{4},\qquad
 \|G_{\rm w}r_{\rm alt}\|_\infty=\frac{21}{88}.
\]
同向失衡与交错失衡的最大温度误差相差整整 $22$ 倍。只取 $C_n^*\|r\|_\infty$ 会给两者一模一样的答案。

\begin{table}[ht]\centering\small
\caption{四格精确实例。每个单元单独作为一块；端点模板覆盖整个允许的边界参数区间。表中仅为方便阅读而四舍五入。}
\begin{tabular}{lrrrrr}\toprule
残差 & 标量 & 位置 & 正负分离 & 双端点 & 实际误差\\\midrule
四格同向 & 5.250 & 5.250 & 5.250 & 5.250 & 5.250\\
交错正负 & 5.250 & 5.250 & 2.744 & 2.600 & 0.239\\
仅一格非零 & 5.250 & 1.600 & 1.600 & 1.600 & 1.600\\\bottomrule
\end{tabular}
\end{table}

表中第二行说明符号在位置幅度不变时仍有用：两个符号界依次为
$483/176$ 和 $8007/3080$。第三行说明即使符号都非负，利用非零残差的空间支撑也能收紧上界，其格林响应最大值为 $12957/8096$。这里求界必须覆盖所有允许的 $\mu$，实际误差则取该行指定的最弱边界；因而即使块已细至单元，双端点估计对混合符号也可能明显偏大。

\paragraph{只移动一个残差格子。}
在 $3\times3$ 网格上取四边均为 $5$，只让一个格子的残差为 $+1$、其余为零。把这一格从角点移到中心，实际最大误差从 $2252723/7061460\approx0.3190$ 变为 $661/1720\approx0.3843$。同一组残差值、同一边界、同样的正号，仅位置不同就得到不同的答案；精确记录也提供覆盖全部许可参数的逐格模板上界。

\paragraph{一个容易写错的捷径。}
从 $G_\mu\le G_{\rm w}$ 不能直接推出 $\|G_\mu r\|_\infty\le\|G_{\rm w}r\|_\infty$。
仍取四格，令 $r=(-46,46,-46,46)^\mathsf{T}/21$，仅把左边的 $h$ 从 $1/5$ 提为 $5$。直接相乘得到 $\|G_{\rm w}r\|_\infty=1$，实际却为 $437/269>1$。因此，带符号残差须按命题 2 处理，不能照搬对非负向量的单调性。

\paragraph{独立精确核对。}
项目用 Python 标准库分数重新组装四格算子并进行有理数消元，逐格核对区间、界的排序和嵌套分块。测试覆盖 16 个边界角点及两组内部参数、全部 $3^4=81$ 个三值残差，以及 1、2、4 块三种嵌套划分，共 $18\times81=1458$ 个参数与残差组合、17496 个单元与分块核对。有限核对能发现实现的符号或排序错误；全参数结论来自命题 1 和命题 2。
